\chapter{Capstone projects}
\label{ch:capstone}

\begin{quote}
\emph{``The best way to learn programming is to build something you care about.''}
\end{quote}

% ============================================================
\section{Overview}

This chapter presents five capstone projects that integrate concepts from the entire course. Each project requires data handling, computation, visualization, and a written report. Choose one (or more) and work through it over 2--3 weeks.

The projects are ordered roughly by domain:
\begin{itemize}
  \item \textbf{Projects 1 and 4:} Dynamical systems and differential equations
  \item \textbf{Project 2:} Deep learning and computer vision
  \item \textbf{Project 3:} Data science and public health
  \item \textbf{Project 5:} Optimization and finance
\end{itemize}

\begin{warningbox}
These are open-ended projects, not step-by-step tutorials. You are expected to consult documentation, experiment, debug, and make design decisions. The instructions provide a roadmap, not a complete solution.
\end{warningbox}

% ============================================================
\section{Project 1: Lotka--Volterra predator-prey simulation}

The Lotka--Volterra equations model predator-prey dynamics:
\[
\frac{dx}{dt} = \alpha x - \beta xy, \qquad \frac{dy}{dt} = \delta xy - \gamma y
\]
where $x$ is prey population, $y$ is predator population, and $\alpha,\beta,\delta,\gamma > 0$.

\begin{enumerate}
  \item Implement using DifferentialEquations.jl with parameters $\alpha=1.1$, $\beta=0.4$, $\delta=0.1$, $\gamma=0.4$.
  \item Create three plots: population vs.\ time, phase portrait ($x$ vs.\ $y$), direction field.
  \item Vary $\alpha$ from 0.5 to 2.0 and overlay phase portraits.
  \item Add stochastic noise via \texttt{SDEProblem}. Plot 20 trajectories alongside the deterministic solution.
  \item Write a function that computes the oscillation period numerically.
\end{enumerate}

\begin{minted}{julia}
using DifferentialEquations, Plots

function lotka_volterra!(du, u, p, t)
    x, y = u; α, β, δ, γ = p
    du[1] = α*x - β*x*y
    du[2] = δ*x*y - γ*y
end

prob = ODEProblem(lotka_volterra!, [10.0,10.0], (0.0,100.0), (1.1,0.4,0.1,0.4))
sol  = solve(prob, Tsit5())
# Your work: create plots, add noise, analyze sensitivity
\end{minted}

% ============================================================
\section{Project 2: Image classifier with Flux.jl}

FashionMNIST contains 70,000 grayscale $28\times28$ images of clothing items in 10 categories: T-shirt, trouser, pullover, dress, coat, sandal, shirt, sneaker, bag, ankle boot.

\begin{enumerate}
  \item Load with MLDatasets.jl. Visualize 16 random samples with labels.
  \item Build a dense network (2+ hidden layers, dropout). Train 15 epochs.
  \item Build a CNN (\texttt{Conv}, \texttt{MaxPool}, \texttt{Dense}). Train 15 epochs.
  \item Plot training loss and test accuracy vs.\ epoch for both architectures.
  \item Compute and display the confusion matrix on the test set.
  \item Identify the 10 most-misclassified images with predicted and true labels.
\end{enumerate}

\begin{minted}{julia}
using Flux, MLDatasets
train_x, train_y = MLDatasets.FashionMNIST(split=:train)[:]

cnn = Chain(
    Conv((3,3), 1=>16, relu, pad=1), MaxPool((2,2)),
    Conv((3,3), 16=>32, relu, pad=1), MaxPool((2,2)),
    Flux.flatten, Dense(32*7*7=>128, relu), Dropout(0.3), Dense(128=>10))
# Your work: train both models, compare, build confusion matrix
\end{minted}

% ============================================================
\section{Project 3: African health data pipeline}

Build a complete data analysis pipeline from public health data sources.

\begin{enumerate}
  \item Download real data from at least two sources: WHO (\url{https://www.who.int/data/gho}), World Bank (\url{https://data.worldbank.org}), Our World in Data (\url{https://ourworldindata.org}).
  \item Clean: handle missing values, inconsistent country names, duplicates.
  \item Join datasets on country and year. Filter for African Union member states.
  \item Compute summary statistics by region (West, East, Central, Southern, North Africa).
  \item Create four publication-quality CairoMakie figures: (a) life expectancy by country, (b) time series for 5 countries, (c) GDP vs.\ health indicator, (d) faceted regional comparison.
  \item Save all figures as PDF.
\end{enumerate}

\begin{minted}{julia}
using CSV, DataFrames, Downloads, CairoMakie, Statistics

url = "https://raw.githubusercontent.com/owid/owid-datasets/" *
      "master/datasets/Life%20expectancy%20-%20OWID/" *
      "Life%20expectancy%20-%20OWID.csv"
Downloads.download(url, "life_exp.csv")
df = CSV.read("life_exp.csv", DataFrame)

african = ["Nigeria","Ethiopia","Egypt","South Africa","Kenya",
    "Tanzania","Ghana","Senegal","Cameroon","Rwanda",
    "Morocco","Tunisia","Algeria","Benin","Togo"]
df_africa = filter(:Entity => in(african), df)
# Your work: clean, join additional sources, analyze, visualize
\end{minted}

\begin{juliatip}
Always use real public data before resorting to synthetic data, especially for African countries. The World Bank, WHO, and Our World in Data provide excellent open datasets with good coverage.
\end{juliatip}

% ============================================================
\section{Project 4: Stochastic SIR epidemic model}

\begin{enumerate}
  \item Implement the deterministic SIR as an \texttt{ODEProblem}.
  \item Add multiplicative noise: $dS = -\beta SI\,dt - \sigma SI\,dW$.
  \item Solve with \texttt{SDEProblem} and \texttt{EM()} (Euler--Maruyama).
  \item Generate 100 trajectories using \texttt{EnsembleProblem}.
  \item Plot the mean with a 95\% confidence band (2.5th and 97.5th percentiles).
  \item Compute the distribution of peak infection time and fraction.
  \item Fit deterministic parameters to the ensemble mean using Optimization.jl.
  \item Investigate the effect of noise intensity $\sigma$ on epidemic outcome.
\end{enumerate}

\begin{minted}{julia}
using DifferentialEquations, Plots, Statistics

function sir_drift!(du, u, p, t)
    S,I,R = u; β,γ,σ = p; N = S+I+R
    du[1] = -β*S*I/N; du[2] = β*S*I/N - γ*I; du[3] = γ*I
end
function sir_noise!(du, u, p, t)
    S,I,R = u; β,γ,σ = p; N = S+I+R
    du[1] = -σ*S*I/N; du[2] = σ*S*I/N; du[3] = 0.0
end

prob = SDEProblem(sir_drift!, sir_noise!, [0.99,0.01,0.0], (0.0,200.0), (0.3,0.1,0.05))
ensemble_sol = solve(EnsembleProblem(prob), EM(); dt=0.1, trajectories=100)
# Your work: plot confidence bands, analyze peak distribution
\end{minted}

% ============================================================
\section{Project 5: Portfolio optimization}

Markowitz optimization: $\min_{\mathbf{w}} \mathbf{w}^\top\Sigma\mathbf{w}$ subject to $\mathbf{w}^\top\boldsymbol{\mu} \geq r_{\min}$, $\mathbf{w}^\top\mathbf{1}=1$, $w_i \geq 0$.

\begin{enumerate}
  \item Generate (or download) daily returns for 10 assets over 5 years.
  \item Compute expected return $\boldsymbol{\mu}$ and covariance $\Sigma$.
  \item Solve the minimum-variance portfolio with Optimization.jl.
  \item Trace the efficient frontier (50 target returns).
  \item Plot the frontier, individual assets, and minimum-variance portfolio.
  \item Add a risk-free asset and compute the tangency portfolio (max Sharpe ratio).
  \item Backtest: compare optimized, equal-weight, and market-cap-weighted portfolios.
\end{enumerate}

\begin{minted}{julia}
using Optimization, OptimizationOptimJL, LinearAlgebra, Plots
using Random; Random.seed!(42)

n_assets = 10
μ = 0.05 .+ 0.15 .* rand(n_assets)
A = randn(n_assets, n_assets)
Σ = (A * A') / 1260 .+ 0.01 * I

portfolio_var(w, p) = w' * p[1] * w   # objective
# Your work: add constraints, trace frontier, backtest
\end{minted}

% ============================================================
\section{Report template}

Your project report should follow this structure:

\begin{enumerate}
  \item \textbf{Title and abstract} (200 words). State the problem, approach, and main result.
  \item \textbf{Introduction} (1 page). Context, motivation, and objectives.
  \item \textbf{Methods} (2--3 pages). Mathematical formulation, algorithms, Julia packages, implementation details.
  \item \textbf{Results} (2--3 pages). Figures, tables, quantitative analysis. Every figure must have a caption.
  \item \textbf{Discussion} (1 page). Interpretation, limitations, possible extensions.
  \item \textbf{Code appendix}. Complete runnable code or a link to a Git repository.
\end{enumerate}

% ============================================================
\section{Grading rubric}

\begin{longtable}{p{4.5cm}p{1.5cm}p{7.5cm}}
\toprule
\textbf{Criterion} & \textbf{Points} & \textbf{Description} \\
\midrule
\endhead
Correctness & 30 & Code runs without errors, results are correct, edge cases handled. \\
\addlinespace
Code quality & 20 & Clean Julia code. Type-stable functions, no globals in hot paths, proper dispatch. \\
\addlinespace
Visualization & 20 & Publication-quality figures with labels, legends, captions. At least 3 plot types. \\
\addlinespace
Report & 15 & Clear writing, logical structure, proper math notation, honest limitations. \\
\addlinespace
Creativity & 15 & Going beyond the minimum: extra analyses, real data, interactive elements. \\
\bottomrule
\end{longtable}

\begin{performancetip}
Use a Git repository from day one. Commit early and often. Include a \texttt{Project.toml} so anyone can reproduce your environment with \texttt{Pkg.activate(".")} and \texttt{Pkg.instantiate()}.
\end{performancetip}

% ============================================================
\section{Presentation guidelines}

Each project will be presented in a 15-minute talk (12 minutes presentation + 3 minutes questions):

\begin{enumerate}
  \item \textbf{Slide 1:} Title, your name, project number.
  \item \textbf{Slides 2--3:} Problem statement and motivation. Why does this problem matter?
  \item \textbf{Slides 4--5:} Methods. What packages did you use? What algorithms? Show key equations.
  \item \textbf{Slides 6--8:} Results. Show your best figures. Explain what each figure reveals.
  \item \textbf{Slide 9:} Discussion. What worked well? What was difficult? What would you do differently?
  \item \textbf{Slide 10:} Conclusion and key takeaways.
\end{enumerate}

\begin{warningbox}
Do not put code on slides unless it illustrates a specific design decision or a clever use of Julia. Slides should show results, figures, and high-level ideas --- not line-by-line code. Keep the code in your report appendix.
\end{warningbox}

\begin{exercise}
\begin{enumerate}
  \item Choose a project. Create a Git repository with \texttt{Project.toml} and an initial script.
  \item Write a one-page project plan: data sources, methods, planned figures, weekly timeline.
  \item Implement a minimal working version that produces at least one correct plot. Commit as ``v0.1.''
  \item Iterate: improve the model, add analyses, refine figures. Commit each improvement.
  \item Write your report and prepare your presentation. Submit the repository, PDF report, and slides.
\end{enumerate}
\end{exercise}

% ============================================================
\section{Chapter summary}

\begin{itemize}
  \item Five capstone projects integrate all course topics: data wrangling, visualization, performance, scientific computing, and machine learning.
  \item Project 1 (Lotka--Volterra) focuses on ODEs, phase portraits, and stochastic simulation.
  \item Project 2 (FashionMNIST) compares dense and convolutional neural networks.
  \item Project 3 (African health data) builds a pipeline from public sources to publication figures.
  \item Project 4 (stochastic SIR) explores ensemble simulation and uncertainty quantification.
  \item Project 5 (portfolio optimization) applies constrained optimization to finance.
  \item Good scientific computing requires clean code, reproducible environments, honest reporting, and clear visualization.
\end{itemize}
